\documentclass[10pt]{amsart}
\usepackage[a4paper,margin=22mm]{geometry}
\usepackage[no-math]{fontspec}
\setmainfont{Latin Modern Roman}[SmallCapsFont={lmromancaps10-regular.otf}]
\usepackage{amsmath,amssymb,amsthm,mathtools,hyperref,amscd,enumitem}
\hypersetup{hidelinks}
\emergencystretch=3em
\newtheorem{theorem}{Theorem}[section]
\newtheorem{lemma}[theorem]{Lemma}
\newtheorem{prop}[theorem]{Proposition}
\newtheorem{coro}[theorem]{Corollary}
\theoremstyle{definition}
\newtheorem{defi}[theorem]{Definition}
\newtheorem{rema}[theorem]{Remark}
\providecommand{\ie}{i.e.}
\providecommand{\C}{\mathcal C}
\DeclareMathOperator\lk{lk}
\DeclareMathOperator\st{st}
\DeclareMathOperator\cost{cost}
\DeclareMathOperator{\Hilb}{Hilb}
\DeclareMathOperator{\Soc}{Soc}
\DeclareMathOperator\depth{depth}
\DeclareMathOperator\Ker{Ker}
\DeclareMathOperator\Coker{Coker}

\newcommand{\complet}[1]{\hat{#1}}
\newcommand{\SB}{\mathcal{SB}}
\newcommand{\field}{\mathbf{k}}
\newcommand{\R}{{\mathbb R}}
\newcommand{\Q}{{\mathbb Q}}
\newcommand{\Z}{{\mathbb Z}}
\newcommand{\ZZ}{{\mathbb Z}}
\renewcommand{\C}{{\mathcal C}}
\newcommand{\K}{{\mathcal K}}
\newcommand{\T}{{\mathcal T}}
\newcommand{\Sm}{{\mathcal S}}
\newcommand{\M}{{\mathcal M}}
\newcommand{\HH}{{\mathcal H}}
\newcommand{\mideal}{\ensuremath{\mathfrak{m}}}



\newcommand{\halffloor}{\left\lfloor \frac{d}{2}\right\rfloor}








%%%%%%%%% a placer avant \begin{document}


%%%%%%%%%%%%%%%%%%%%%%%%%%%%%%%%%%%%%
\newcommand*{\mk}{\mkern -1mu}
\newcommand*{\Mk}{\mkern -2mu}
\newcommand*{\mK}{\mkern 1mu}
\newcommand*{\MK}{\mkern 2mu}

\hypersetup{urlcolor=purple, linkcolor=blue, citecolor=red}


\newcommand*{\romanenumi}{\renewcommand*{\theenumi}{\roman{enumi}}}
\newcommand*{\Romanenumi}{\renewcommand*{\theenumi}{\Roman{enumi}}}
\newcommand*{\alphenumi}{\renewcommand*{\theenumi}{\alph{enumi}}}
\newcommand*{\Alphenumi}{\renewcommand*{\theenumi}{\Alph{enumi}}}
%%%%%%%%%%%%%%%%%%%%%%%%%%%%%%%%%%



%%%%%%%%%%%%%%%%%%%%%%%%%%%%%%%%%%%%%%%%%%%
\title[Socles and Hilbert numbers for completed manifolds]{Socle dimensions and corrected Hilbert numbers of completed homology manifolds with boundary}
\author{Isabella Novik and Ed Swartz}
\date{Algebraic Combinatorics3(4)(2020),887–911; DOI10.5802/alco.121}
\begin{document}
\maketitle
\setcounter{section}{1}
\section{Preliminaries} \label{sec:Prelim}
In this section we review the necessary background material on simplicial complexes and their Stanley--Reisner rings with a special emphasis on homology manifolds with and without boundary as well as on singular simplicial complexes that have only one singular vertex. We refer the reader to~\cite[Chapter~2]{Stanley-96} and the papers~\cite{Novik-Swartz-09:Socles, Novik-Swartz-12} for more details on the subject. 

\subsection{Simplicial complexes: homology manifolds and their completions}
A \emph{simplicial complex} $\Delta$ on the ground set $V$ is a collection of subsets of $V$ that is closed under inclusion.
The maximal faces (with respect to inclusion) are called \emph{facets}. The \emph{dimension of a face} $F\in \Delta$ is $\dim F:=|F|-1$ and
the \emph{dimension of $\Delta$} is the maximal dimension of its faces. A complex is \emph{pure} if all of its facets have the same dimension. A complex $\Delta$ is \emph{$j$-neighborly} if every $j$-element subset of $V$ is a face of $\Delta$. 

Let $\Delta$ be a simplicial complex and let $F$ be a face of $\Delta$. The \emph{star} and the \emph{link} of $F$ in $\Delta$ are the following subcomplexes of~$\Delta$: 
\[
\st F=\st_\Delta F:=\{G\in \Delta \, \mid \, G\cup F\in \Delta\}, \quad
 \lk F=\lk_\Delta F:=\{G\in \st_\Delta F \, \mid \, G\cap F=\emptyset\}.
\]
In particular, the link of the empty face is the complex $\Delta$ itself. We refer to the links of non-empty faces as \emph{proper links}. The \emph{contrastar} of $F$ in $\Delta$ (also known as the \emph{deletion} of $F$ from $\Delta$) is $\cost F=\cost_\Delta F:=\{G\in \Delta \mid G\not\supseteq F\}$. If $v$ is a vertex (\ie a $0$-dimensional face), it is customary to write $v\in\Delta$, $\st v$, $\lk v$, and $\cost v$ instead of $\{v\}\in\Delta$, $\st \{v\}$, $\lk\{v\}$, and $\cost\{v\}$. (In fact, we will sometimes write $\Delta-v$ instead of $\cost\{v\}$.)
Also, if $v_0\notin V$ is a new vertex, then the \emph{cone} over $\Delta$ with apex $v_0$ is $v_0\ast\Delta:=\Delta \cup \{v_0\cup F \mid F\in \Delta\}$.


Throughout the paper, we fix an \emph{infinite} field $\field$. We denote by $\tilde{H}_\ast(\Delta;\field)$ the \emph{reduced simplicial homology} of $\Delta$ with coefficients in $\field$ and by $\tilde{\beta}_i(\Delta):=\dim_\field \tilde{H}_i(\Delta;\field)$ the $i$-th \emph{reduced Betti number}. Occasionally, we also use the (reduced) relative simplicial homology of a pair $(\Delta,\Gamma)$: $\tilde{H}_i(\Delta, \Gamma;\field)$ and the corresponding Betti numbers $\tilde{\beta}_i(\Delta, \Gamma):=\dim_\field \tilde{H}_i(\Delta, \Gamma;\field)$. We remark that $\tilde{H}_i(\Delta, \Gamma;\field)=H_i(\Delta, \Gamma;\field)$ for all $\Gamma\subseteq \Delta$ and all $i>0$ and that $\tilde{H}_0(\Delta, \Gamma;\field)=H_0(\Delta, \Gamma;\field)$ if $\Gamma\neq \emptyset$. On the other hand, $\tilde{\beta}_0(\Delta,\emptyset)=\tilde{\beta}_0(\Delta)=\beta_0(\Delta,\emptyset)-1$.

One of the central objects of this paper is a \emph{$\field$-homology manifold}. A pure $(d-1)$-dimensional simplicial complex $\Delta$ is a \emph{$\field$-homology manifold without boundary} (or a \emph{closed $\field$-homology manifold}) if the homology (computed over $\field$) of every proper link of $\Delta$, $\lk_\Delta F$, coincides with the homology of a $(d-1-|F|)$-dimensional sphere. In this case, we write $\partial \Delta = \emptyset$. Similarly, a pure $(d-1)$-dimensional simplicial complex $\Delta$ is a \emph{$\field$-homology manifold with boundary} if every proper link of $\Delta$, $\lk_\Delta F$, has the homology of a $(d-1-|F|)$-dimensional sphere or a ball (over $\field$), and the \emph{boundary complex} of $\Delta$, \ie  
\[ 
\partial \Delta:= \left\{F \in \Delta \mid \tilde{H}_\ast(\lk_\Delta F; \field)=0\right\}\cup\{\emptyset\}
\]
is a $(d-2)$-dimensional $\field$-homology manifold without boundary. The faces of $\partial\Delta$ are called \emph{boundary faces}. The non-boundary faces of $\Delta$ are called \emph{interior faces}. When the field plays no role we simply call $\Delta$ a homology manifold (with or without boundary). We refer the reader to Chapter 8 of Munkers' book~\cite{Munkres-book} (and especially \S\S~63, 65, 70 and 72 there) for more background on homology manifolds.

The prototypical example of a homology manifold (with or without boundary) is a triangulation of a topological manifold (with or without boundary). A connected $\field$-homology manifold $\Delta$ is \emph{orientable} if the top homology of the pair $(\Delta, \partial \Delta)$ is $1$-dimensional. In this case, $(\Delta, \partial \Delta)$ satisfies the usual Poincar{\'e}--Lefschetz duality associated with orientable compact manifolds. Note that an arbitrary triangulation of any topological manifold (orientable or not, with or without boundary) is an \emph{orientable} $\Z/2\Z$-homology manifold.

A $\field$-homology $(d-1)$-sphere is a $(d-1)$-dimensional $\field$-homology manifold without boundary that has the same homology as the $(d-1)$-dimensional sphere. A $\field$-homology $(d-1)$-ball is a $(d-1)$-dimensional $\field$-homology manifold with boundary whose homology is trivial and whose boundary complex is a $\field$-homology $(d-2)$-sphere. The contrastar of any vertex in a $\field$-homology $(d-1)$-sphere is a $\field$-homology $(d-1)$-ball. Furthermore, every proper link of a $\field$-homology manifold without boundary is a $\field$-homology sphere, while a proper link of a $\field$-homology manifold with boundary is a $\field$-homology sphere or ball.

Let $\Delta$ be a $\field$-homology manifold with or without boundary and let $v_0\not\in V$ be a new vertex. A key to most of our proofs is the \emph{completion of $\Delta$}, $\complet{\Delta}$, defined as follows:
\[
\complet{\Delta}:=\Delta\cup (v_0\ast\partial\Delta).
\]
 Note that we define $v_0\ast\emptyset =\emptyset$; hence if $\Delta$ is a homology manifold without boundary, then $\complet{\Delta}=\Delta$. 

A pure simplicial complex $\Gamma$ is a \emph{complex with at most one singularity} if all of the vertex links of $\Gamma$ but possibly one are $\field$-homology balls or spheres. This exceptional vertex is called a \emph{singular} vertex; the other vertices are called \emph{non-singular}. For instance, if $\Delta$ is a $\field$-homology manifold with boundary, $\complet{\Delta}$ is a completion of $\Delta$, and $v\neq v_0$, then both $\complet{\Delta}$ and $\cost_{\complet{\Delta}}v$ are complexes with (at most) one singular vertex, namely, $v_0$.

When only topological properties of a space are relevant we may use capital roman letters. For instance, ``If $X$ is a $d$-dimensional ball, then its boundary $Y$ is a $(d-1)$-dimensional sphere.''

%%%%%%%%%%%%%%%%%%%%%%%%%%%%%%%%%%

\subsection{Face numbers and the Stanley--Reisner rings} Let $\Delta$ be a $(d-1)$-dimensional simplicial complex on the vertex set $V$. Denote by $f_i(\Delta)$ the number of $i$-dimensional faces of $\Delta$.
The \emph{$f$-vector} of $\Delta$ is $f(\Delta)=(f_{-1}(\Delta), f_0(\Delta),\ldots, f_{d-1}(\Delta))$ and the \emph{$h$-vector} of $\Delta$ is $h(\Delta)=(h_0(\Delta), h_1(\Delta), \ldots, h_d(\Delta))$, where 
\[h_i(\Delta):=\sum_{j=0}^i (-1)^{i-j}\binom{d-j}{d-i}f_{j-1}(\Delta).\]

Let $A=\field[x_v \mid v\in V]$ be a polynomial ring, and let $\mideal=(x_v \mid v\in V)$ be the graded maximal ideal of $A$. For $F\subseteq V$, write $x_F:=\prod_{v\in F} x_v$. The \emph{Stanley--Reisner ideal} $I_\Delta$ of $\Delta$ is the ideal of $A$ defined by 
\[
I_\Delta=(x_F \mid F\subseteq V, F\notin\Delta).\]
The \emph{Stanley--Reisner ring} $\field[\Delta]$ of $\Delta$ (over $\field$) is the quotient ring $\field[\Delta]=A/I_\Delta$. In particular, $\field[\Delta]$ is a graded ring; it is also a graded $A$-module. If $\dim\Delta=d-1$, then the Krull dimension of $\field[\Delta]$ is $d$ and the Hilbert series of $\field[\Delta]$ is given by (\cite[Chapter~II.1]{Stanley-96})
\[
\Hilb(\field[\Delta]; \lambda)=\frac{\sum_{i=0}^d h_i(\Delta)\lambda^i}{(1-\lambda)^d}.\]

A \emph{linear system of parameters} (or \emph{l.s.o.p} for short) for $\field[\Delta]$ is a sequence $\Theta=\theta_1,\ldots,\theta_d$ of $d=\dim\Delta + 1$ linear forms in $\mideal$ such that \[\field(\Delta,\Theta):= \field[\Delta]/\Theta\field[\Delta]\]
has Krull dimension zero (\ie it is a finite-dimensional $\field$-space). Since $\field$ is infinite, by the Noether normalization lemma, an l.s.o.p.~always exists: a generic choice of $\theta_1,\ldots,\theta_d$ does the job.
The ring $\field(\Delta,\Theta)$ is called an \emph{Artinian reduction} of $\field[\Delta]$. 

We need a few more definitions. If $M$ is a finitely-generated graded $A$-module, we let $M_j$ denote the $j$-th homogeneous component of $M$. For $\tau\in A$, define $(0:_M \tau):=\{\nu\in M \mid \tau\nu=0\}$. The \emph{socle} of $M$ is the following graded submodule of $M$:
\[\Soc M=\bigcap_{v\in V}(0:_M x_v)=\{\nu\in M \mid \mideal\nu=0\}.\]
In particular, for any choice of integers $i_1<i_2<\cdots<i_\ell$, $\bigoplus_{j=1}^\ell (\Soc M)_{i_j}$ is a submodule of $M$.
For a standard graded $\field$-algebra $M=A/I$ of Krull dimension zero, this allows us to define the \emph{interior socle} of $M$: 
\[
\Soc^\circ M:=\bigoplus_{i=0}^{d_0-1}(\Soc M)_i, \quad \mbox{where $d_0:=\max\{j \mid M_j\neq 0\}$}.
\]
\looseness-1
We say that $A/I$ is a \emph{level algebra} if $\Soc^\circ(A/I)=0$, and that $A/I$ is \emph{Gorenstein} if it has a $1$-dimensional socle. Equivalently, $A/I$ is Gorenstein if it is level and $\dim_\field (A/I)_{d_0}=1$.


We are interested in the Hilbert functions of $\field(\Delta,\Theta)$ and its quotient \[\overline{\field(\Delta,\Theta)}:= \field(\Delta,\Theta)/\Soc^\circ \field(\Delta,\Theta).\] 
\begin{defi} \label{def-h'-h''}
Let $\Delta$ be a $(d-1)$-dimensional simplicial complex and let $\Theta=\theta_1,\ldots,\theta_d$ be a \emph{generic} l.s.o.p.~for $\field[\Delta]$. The \emph{$h'$-} and the \emph{$h''$-numbers} of $\Delta$ are defined by
\[
h'_j(\Delta):=\dim_\field \field(\Delta,\Theta)_j \quad \mbox{and} \quad h''_j(\Delta):=\dim_\field \overline{\field(\Delta,\Theta)}_j \quad \mbox{(for $j\geq 0$), respectively.}
\]
\end{defi} 

Although $\field$ is suppressed from our notation, the $h'$- and $h''$-numbers do depend on $\field$. For any $(d-1)$-dimensional simplicial complex $\Delta$, $h'_j(\Delta)=0$ for all $j>d$ (see~\cite[Proposition~III.2.4(b)]{Stanley-96}), while $h'_d(\Delta)=h''_d(\Delta)=\tilde{\beta}_{d-1}(\Delta)$ (see~\cite[Theorem~4.1]{TayWhiteley} and~\cite[Lemma~2.2(3)]{Babson-Novik}). The following theorem collects several other known results on $h'$- and $h''$-numbers.

\goodbreak
\begin{theorem} \label{known-h'-h''} Let $\Delta$ be a $(d-1)$-dimensional simplicial complex.
\begin{enumerate}
\item\label{theo2.2_1} If $\Delta$ is a $\field$-homology sphere or a ball, then $h'_i(\Delta)=h_i(\Delta)$ for all $0\leq i\leq d$.
\item\label{theo2.2_2} If $\Delta$ is a $\field$-homology manifold with or without boundary, then 
\[h'_i(\Delta)=h_i(\Delta)-\binom{d}{i}\sum_{j=1}^{i-1}(-1)^{i-j}\tilde{\beta}_{j-1}(\Delta) \qquad \forall ~0\leq i\leq d.\]
\item\label{theo2.2_3} If $\Delta$ is a connected, orientable $\field$-homology manifold without boundary, then 
\[h''_i(\Delta)=h'_i(\Delta)-\binom{d}{i}\tilde{\beta}_{i-1}(\Delta)=
 h_i(\Delta)-\binom{d}{i}\sum_{j=1}^{i}(-1)^{i-j}\tilde{\beta}_{j-1}(\Delta) \qquad \forall ~ 0\leq i\leq d-1.\]
\item\label{theo2.2_4} If $\Delta$ is a complex with (at most) one singular vertex $u$, then for all $0\leq i\leq d$,
\[h'_i(\Delta)=h_i(\Delta)-\sum_{j=1}^{i-1}(-1)^{i-j}\left(\binom{d-1}{i-1}\tilde{\beta}_{j-1}(\Delta)+\binom{d-1}{i}\tilde{\beta}_{j-1}(\cost_{\Delta}u)\right).\]
\end{enumerate}
\end{theorem}

\noindent Part 1 of this theorem is due to Stanley~\cite{Stanley-75}, part 2 is due to Schenzel~\cite{Schenzel-81}, part 3 is~\cite[Theorem~1.3]{Novik-Swartz-09:Gorenstein}, and part 4 is a special case of~\cite[Theorem~4.7]{Novik-Swartz-12}. When $\Delta$ is a $\field$-homology manifold with boundary, part 4 allows us to compute the $h'$-numbers of $\complet{\Delta}$. One of the goals of this paper is to understand the $h''$-numbers of $\complet{\Delta}$ where we in addition assume that $\Delta$ is connected and orientable. This requires some results on the local cohomology of $\field[\complet{\Delta}]$ that we review in the next subsection.


 It is worth pointing out that there are several reasons for working with $\complet{\Delta}$ instead of $\Delta$ itself. One of the reasons is that, as we will see in Section 3, the ring $\overline{\field(\complet{\Delta},\Theta)}$ is Gorenstein. Another reason is that the $h''$-numbers of $\field$-homology manifolds with boundary are hard to calculate. In view of Schenzel's formula (Theorem~\ref{known-h'-h''}\eqref{theo2.2_2}), to compute the $h''$-numbers of a homology manifold $\Delta$ (with or without boundary), one needs to understand the module $\Soc \field(\Delta,\Theta)$. Theorems 2.2 and 3.4 in~\cite{Novik-Swartz-09:Socles} decompose this module into two parts: the first part involves the well-understood local cohomology modules of $\field[\Delta]$ while the second part involves a certain mysterious submodule of $\Soc H^d(\field[\Delta])$. If $\Delta$ is a connected, orientable $\field$-homology manifold without boundary, then, as was shown in~\cite[Theorem~2.1]{Novik-Swartz-09:Gorenstein}, the socle of $H^d(\field[\Delta])$, and hence also the ``mysterious submodule'', vanish in all but the $0$-th degree; this leads to the proof of Theorem~\ref{known-h'-h''}\eqref{theo2.2_3}. However, for $\field$-homology manifolds with boundary, at present we are lacking even a conjectural description of this mysterious part. For instance, if $\Delta$ is a $2$-dimensional disk then the $h$-vector of $\Delta$ is $(1,m,n,0)$ with $m \ge n \ge 0.$ If $n >0,$ then $0 \le \dim_\field \Soc^\circ \field(\Delta,\Theta) \le m-n$ and every value in the inequality is possible. In particular, the topology of $\Delta$ does not determine $\dim_\field \Soc^\circ \field (\Delta, \Theta).$


\subsection{Local cohomology and Gr{\"a}be's theorem}
Let $M$ be an arbitrary finitely-generated graded $A$-module. We denote by $H^i_\mideal(M)$ the $i$-th local cohomology of $M$ with respect to $\mideal$. 

For a simplicial complex $\Delta$, Gr{\"a}be~\cite{Grabe-84} gave a description of $H^i_\mideal(\field[\Delta])$ and its $A$-module structure in terms of simplicial cohomology of the links of $\Delta$ and the maps between them. When $\Delta$ is a complex with one singular vertex $u$, this description takes the following simple form. For $F\in\Delta$, consider the $i$-th simplicial cohomology of the pair $(\Delta, \cost_\Delta F)$ with coefficients in $\field$:
\[
H^i_F(\Delta):=\tilde{H}^i(\Delta, \cost_\Delta F; \field)\cong \tilde{H}^{i-|F|}(\lk_\Delta F; \field). \]
In particular, $H^i_\emptyset\Mk =\Mk \tilde{H}^i(\Delta, \emptyset;\field)\Mk =\Mk \tilde{H}^i(\Delta; \field)$. If $G\Mk \subseteq \Mk F\Mk \in\Mk \Delta$, we let $\iota^\ast\Mk :\Mk  H^i_F(\Delta)\Mk \to\Mk  H^i_G(\Delta)$ be the map induced by inclusion $\iota: \cost_\Delta G \to \cost_\Delta F$.

\begin{theorem}[Gr{\"a}be] \label{local-cohom-1-singul}
%{\rm[Gr{\"a}be]} 
Let $\Delta$ be a $(d-1)$-dimensional simplicial complex with one singular vertex $u$, and let $-1\leq i <d-1$. Then
\[
H^{i+1}_\mideal(\field[\Delta])_{-j}=\begin{cases}
0 & \mbox{ if $j<0$},\\
H^i_\emptyset(\Delta) & \mbox{ if $j=0$},\\
H^i_{\{u\}}(\Delta) & \mbox{ if $j>0$}.
\end{cases}
\]
For every vertex $w\neq u$ and any integer $j$, the map $\cdot x_w: H^{i+1}_\mideal(\field[\Delta])_{-(j+1)} \to H^{i+1}_\mideal(\field[\Delta])_{-j}$ is the zero map; on the other hand, 
\[\cdot x_u= \begin{cases}
\mbox{$0$-map} & \mbox{ if $j<0$},\\
\iota^*: H^i_{\{u\}}(\Delta) \to H^i_\emptyset(\Delta) & \mbox{ if $j=0$},\\
\mbox{identity map: $H^i_{\{u\}}(\Delta) \to H^i_{\{u\}}(\Delta)$} & \mbox{ if $j>0$}.
\end{cases}
\]
\end{theorem}

The description of $H^d_\mideal(\field[\Delta])$ is quite a bit more involved. To this end, for a monomial $\rho\in A$, define the \emph{support} of $\rho$ by $s(\rho):=\{v\in V \mid x_v \mbox{ divides }\rho\}$. Let $\M(\Delta)$ be the set of all monomials in $A$ whose support is in $\Delta$, and let $\M_j(\Delta):=\{\rho\in\M(\Delta) \mid \deg(\rho)=j\}$.

\begin{theorem}[Gr{\"a}be] \label{Grabe-d}
%{\rm[Gr{\"a}be]} 
Let $\Delta$ be any $(d-1)$-dimensional simplicial complex. Then for $j\in\ZZ$, 
\[ H^d_{\mideal}(\field[\Delta])_{-j}=\bigoplus_{\rho\in\M_j(\Delta)} \HH_{\rho}, \quad \mbox{where} \quad \HH_{\rho}=H^{d-1}_{s(\rho)}(\Delta),\]
and the $A$-structure on the $\rho$-th component of the right-hand side is given by
\[
\cdot x_\ell = \begin{cases} 
\mbox{$0$-map} & \mbox{if $\ell\notin s(\rho)$},\\
\iota^*: H^{d-1}_{s(\rho)}(\Delta) \to H^{d-1}_{s(\rho/x_\ell)}(\Delta) & \mbox{if $\ell\in s(\rho)$, but $\ell\notin s(\rho/x_\ell)$},\\
\mbox{identity map: $H^{d-1}_{s(\rho)}(\Delta) \to H^{d-1}_{s(\rho/x_\ell)}(\Delta)$} & \mbox{if $\ell\in s(\rho)$, and $\ell\in s(\rho/x_\ell)$}.
\end{cases}
\]
\end{theorem}

\section{The Gorenstein and Weak Lefschetz properties} \label{sec:Gorenstein}
If $\Delta$ is a $\field$-homology sphere and $\Theta$ is an arbitrary l.s.o.p.~for $\field[\Delta]$, then $\field(\Delta,\Theta)$ is Gorenstein (see~\cite[Theorem~II.5.1]{Stanley-96}). This result was extended in~\cite[Theorem~1.4]{Novik-Swartz-09:Gorenstein} to connected, orientable $\field$-homology manifolds without boundary: if $\Delta$ is such a complex and $\Theta$ is an l.s.o.p.~for $\field[\Delta]$, then $\overline{\field(\Delta,\Theta)}$ is Gorenstein. Here we further extend this result to manifolds with boundary. 

Throughout this section, we let $\Delta$ be a $\field$-homology manifold with boundary. We assume that $\Delta$ is $(d-1)$-dimensional and has vertex set $V$, and so $\complet{\Delta}$ has vertex set $V_0:=V\cup\{v_0\}$. The main result of this section is:

\begin{theorem} \label{thm-Gorenstein}
Let $\Delta$ be a connected, orientable $\field$-homology manifold with boundary and let $\Theta$ be a generic l.s.o.p.~for $\field[\complet{\Delta}]$. Then $\overline{\field(\complet{\Delta},\Theta)}$ is Gorenstein.
\end{theorem}

The proof relies on a few lemmas. For these lemmas we fix a vertex $v$ of $\Delta$. (Hence $v$ is a non-singular vertex of $\complet{\Delta}$.)

\begin{lemma} \label{lem:homol-of-complex-and-costar}
Let $\Delta$ be a $(d-1)$-dimensional, $\field$-homology manifold with boundary, and let $v$ be a vertex of $\Delta$. Then for all $j\leq d-3$,
\[
\tilde{\beta}_j(\complet{\Delta})=\tilde{\beta}_j(\cost_{\complet{\Delta}}v) \quad \mbox{and} \quad \tilde{\beta}_j(\Delta)=\tilde{\beta}_j(\cost_{\Delta}v).
\]
\end{lemma}
\begin{proof}
\looseness-1
The proof is a simple application of the Mayer-Vietoris argument. Indeed, since $v\neq v_0$, the link $\complet{L}:=\lk_{\complet{\Delta}}v$ is a $\field$-homology $(d-2)$-sphere, while the link $L:=\lk_\Delta v$ is a $\field$-homology $(d-2)$-sphere or a $(d-2)$-ball. In either case, $\tilde{\beta}_j(\complet{L})=\tilde{\beta}_j(L)=0$ for all $j\leq d-3$. Since, the stars $\st_\Delta v$ and $\st_{\complet{\Delta}}v$ are acyclic, considering the following portion 
\[
\postdisplaypenalty1000000
\tilde{H}_j(L)\to\tilde{H}_j(\cost_\Delta v)\oplus\tilde{H}_j(\st_\Delta v)\to\tilde{H}_j(\Delta)\to\tilde{H}_{j-1}(L)
\]
of the Mayer-Vietoris sequence for $\Delta$ and its analog for $\complet{\Delta}$ yields the result.
\end{proof}

 The following lemma is a generalization of~\cite[Proposition~4.24]{Swartz-09}. We set $A:=\field[x_u \mid u\in V_0]$ and $A^v:=\field[x_u \mid u\in V_0\setminus\{v\}]$.
Observe that $\field[\complet{\Delta}]$ and $\field[\cost_{\complet{\Delta}}v]$ have natural $A$-module structures (where multiplication by $x_v$ on $\field[\cost_{\complet{\Delta}}v]$ is the zero map), while $\field[\lk_{\complet{\Delta}}v]$ has a natural $A^v$-module structure (if $u\neq v$ is not in the link of $v$, then multiplication by $x_u$ is the zero map). Let $\Theta=\theta_1,\ldots,\theta_d\in A$ be a generic l.s.o.p.~for $\field[\complet{\Delta}]$, and hence also for $\field[\cost_{\complet{\Delta}}v]$. Since $\Theta$ is generic, $\theta_1$ has non-vanishing coefficients. So by scaling the variables if necessary, we can work in an isomorphic setting and assume w.l.o.g.~that \emph{all} coefficients of $\theta_1$ are equal to $1$. 
Let $\theta'_1:=\theta_1-x_v$, and for $j>1$, let $\theta'_j=\theta_j-c_j\theta_1$ where $c_j$ is the coefficient of $x_v$ in $\theta_j$. Then $\theta'_1, \theta'_2,\ldots,\theta'_d$ can be viewed as elements of $A^v$, with $\Theta^v=\{\theta'_2,\ldots,\theta'_d\}\subset A^v$ forming an l.s.o.p.~for $\field[\lk_{\complet{\Delta}}v]$. Furthermore, the ring $\field(\lk_{\complet{\Delta}}v, \Theta^v)$ inherits an $A^v$-module structure, and defining
\[
x_v\cdot y:=-\theta'_1\cdot y \quad \mbox{for }y\in \field(\lk_{\complet{\Delta}}v, \Theta^v)
\]
extends it to an $A$-module structure.

\begin{lemma} \label{analog-of-Prop4.24}
 Let $\Delta$ be a $(d-1)$-dimensional, connected, orientable $\field$-homology manifold with boundary, and let $v$ be a vertex of $\Delta$. Then 
the map
\[\phi_v: \field(\lk_{\complet{\Delta}}v, \Theta^v) \to (x_v)\field(\complet{\Delta},\Theta) \quad \mbox{given by } z \to x_v\cdot z
\]
is well-defined and is a graded isomorphism of $A$-modules (of degree $1$).
\end{lemma}
\begin{proof} The proof of~\cite[Proposition~4.24]{Swartz-09} shows that $\phi_v$ is a well-defined and surjective homomorphism of $A$-modules. Thus to complete the proof, it suffices to check that for $1\leq i\leq d$, the dimensions of $\field$-spaces $\big(\field(\lk_{\complet{\Delta}}v, \Theta^v)\big)_{i-1}$ and $\big((x_v)\field(\complet{\Delta},\Theta)\big)_i$ agree.
Since $v\neq v_0$, the link $\lk_{\complet{\Delta}}v$ is a $\field$-homology sphere, and so 
\begin{equation} \label{dim-lk}
\dim_\field \big(\field(\lk_{\complet{\Delta}}v, \Theta^v)\big)_{i-1}=h_{i-1}(\lk_{\complet{\Delta}}v) \quad \mbox{for all }i\leq d.
\end{equation}

To compute $\dim_\field \big((x_v)\field(\complet{\Delta},\Theta)\big)_i$ for $i\leq d$, consider the following exact sequence, induced by the natural surjection $\field[\complet{\Delta}]\to \field[\cost_{\complet{\Delta}}v]$, 
\begin{equation} \label{exact-x_v}
0\to (x_v)\field(\complet{\Delta},\Theta) \to \field(\complet{\Delta},\Theta) \to \field(\cost_{\complet{\Delta}}v, \Theta)\to 0.
\end{equation}
If $i=d$, then, since $\Delta$ is connected and orientable, 
\[
\dim_\field \big(\field(\complet{\Delta},\Theta)\big)_d=\tilde{\beta}_{d-1}(\complet{\Delta})=1 \mbox{ while } \dim_\field \field(\cost_{\complet{\Delta}}v, \Theta)_d=\tilde{\beta}_{d-1}(\cost_{\complet{\Delta}}v)=0.
\] 
Hence, in this case eq.~\eqref{exact-x_v} implies that
\[
\dim_\field \big((x_v)\field(\complet{\Delta},\Theta)\big)_d \Mk =\Mk \dim_\field \big(\field(\complet{\Delta},\Theta)\big)_d\Mk =\Mk 1
\Mk =\Mk \tilde{\beta}_{d-1}(\lk_{\complet{\Delta}}v)\Mk =\Mk \dim_\field \big(\field(\lk_{\complet{\Delta}}v, \Theta^v)\big)_{d-1},
\]
as desired. 

Thus for the rest of the proof we assume that $1\leq i\leq d-1$. Since both $\complet{\Delta}$ and $\cost_{\complet{\Delta}}v$ are complexes with at most one singular vertex, namely $v_0$, and since $\cost_{\complet{\Delta}}v_0=\Delta$, we infer from Theorem~\ref{known-h'-h''}\eqref{theo2.2_4} that
\begin{multline} \label{h'}
\dim_\field \big(\field(\complet{\Delta},\Theta)\big)_i-h_i(\complet{\Delta}) \\
=
-\sum_{j=1}^{i-1}(-1)^{i-j}\left(\binom{d-1}{i-1}\tilde{\beta}_{j-1}(\complet{\Delta})+\binom{d-1}{i}\tilde{\beta}_{j-1}(\Delta)\right),
\end{multline}
and that a similar expression holds for $\dim_\field\big(\field(\cost_{\complet{\Delta}}v,\Theta)\big)_i-h_i(\cost_{\complet{\Delta}}v)$: to obtain it, simply replace all occurrences of $\complet{\Delta}$ on the right-hand side of~\eqref{h'} with $\cost_{\complet{\Delta}}v$ and those of $\Delta$ with $\cost_{\Delta}v$. Since according to Lemma~\ref{lem:homol-of-complex-and-costar}, for $i\leq d-1$, these replacements do not affect the value of the right-hand side of~\eqref{h'}, we conclude that for all $1\leq i\leq d-1$,
\begin{align*}
\dim_\field \big((x_v)\field(\complet{\Delta},\Theta)\big)_i & 
\ \ \mathop{=}^{\hbox to 0pt {\hss\text{\scriptsize by~\eqref{exact-x_v}}\hss}} \ \ 
%\stackrel{\mbox{\tiny{by~\eqref{exact-x_v}}}}{=} 
\dim_\field\big(\field(\complet{\Delta},\Theta)\big)_i- \dim_\field\big(\field(\cost_{\complet{\Delta}}v,\Theta)\big)_i 
\\
&\ \ = \ \ h_i(\complet{\Delta})-h_i(\cost_{\complet{\Delta}}v) \\
&\ \ = \ \ h_{i-1}(\lk_{\complet{\Delta}}v) \stackrel{\mbox{\tiny{by~\eqref{dim-lk}}}}{=}
\dim_\field \big(\field(\lk_{\complet{\Delta}}v, \Theta^v)\big)_{i-1}, 
\end{align*}
where the penultimate step uses~\cite[Lemma~4.1]{Athanasiadis-11}. The result follows.
\end{proof}

We are now in a position to prove Theorem~\ref{thm-Gorenstein}. Our proof follows the same outline as the proof of~\cite[Theorem~1.4]{Novik-Swartz-09:Gorenstein} with an additional twist at the end.

\begin{proof}[Proof of Theorem~\ref{thm-Gorenstein}]
Let $\Delta$ be a $(d-1)$-dimensional, connected, orientable $\field$-homology manifold with boundary, and let $\Theta$ be a generic l.s.o.p.~for $\field[\complet{\Delta}]$. (As before, we assume w.l.o.g.~that all coefficients of $\theta_1$ are equal to $1$.) Then $\dim_\field \overline{\field(\complet{\Delta},\Theta)}_d =\tilde{\beta}_{d-1}(\complet{\Delta})=1$. Hence, we only need to verify that the socle of $\overline{\field(\complet{\Delta},\Theta)}=\field(\complet{\Delta},\Theta)/\Soc^\circ \field(\complet{\Delta},\Theta)$ vanishes in all degrees $j\neq d$. Since $\big(\Soc^\circ \field(\complet{\Delta},\Theta)\big)_d=0$ and $\big(\Soc^\circ \field(\complet{\Delta},\Theta)\big)_{d-1}=\big(\Soc \field(\complet{\Delta},\Theta)\big)_{d-1}$, this does hold for $j=d-1$.

Now, let $j\leq d-2$, and let $y\in \field(\complet{\Delta},\Theta)_j$ be such that $x_v\cdot y\in (\Soc \field(\complet{\Delta},\Theta))_{j+1}$ for every vertex $v$ of $\complet{\Delta}$. We must show that $y\in \Soc \field(\complet{\Delta},\Theta)$. Assume first that $v\neq v_0$. Then the isomorphism of Lemma~\ref{analog-of-Prop4.24} implies that $y^v:=\phi_v^{-1}(x_v\cdot y)\in (\field(\lk_{\complet{\Delta}}v, \Theta^v))_j$ is in the socle of $\field(\lk_{\complet{\Delta}}v, \Theta^v)$. Since $\lk_{\complet{\Delta}}v$ is a $\field$-homology $(d-2)$-sphere, $\field(\lk_{\complet{\Delta}}v, \Theta^v)$ is Gorenstein, and hence its socle vanishes in all degrees $\leq d-2$. Therefore, $y^v=0$. We conclude that
\begin{equation} \label{v-neq-v0}
x_v\cdot y=\phi_v(y^v)=0 \mbox{ in } \field(\complet{\Delta},\Theta) \quad \mbox{for all } v\neq v_0.
\end{equation}

\looseness-1
Finally, to show that $x_{v_0}\cdot y=0$ in $\field(\complet{\Delta},\Theta)$, recall that $\theta_1=x_{v_0}+\sum_{v\neq v_0}x_v$, and so
\begin{equation} \label{v0}
\theta_1\cdot y=x_{v_0}\cdot y + \sum_{v\neq v_0} x_v\cdot y.
\end{equation}
The left-hand side of~\eqref{v0} is zero in $\field(\complet{\Delta},\Theta)=\field[\complet{\Delta}]/\Theta\field[\complet{\Delta}]$. Furthermore, by~\eqref{v-neq-v0}, all summands on the right-hand side of~\eqref{v0}, except possibly $x_{v_0}\cdot y$, are zeros in $\field(\complet{\Delta},\Theta)$. Thus $x_{v_0}\cdot y$ must be zero in $\field(\complet{\Delta},\Theta)$. The result follows. 
\end{proof}

We now turn to some consequences of Theorem~\ref{thm-Gorenstein}. As the Hilbert function of a Gorenstein graded $\field$-algebra of Krull dimension zero is always symmetric, one immediate corollary is
\begin{coro} \label{symmetry}
Let $\Delta$ be a $(d-1)$-dimensional, connected, orientable $\field$-homology manifold with boundary. Then 
$h''_i(\complet{\Delta})=h''_{d-i}(\complet{\Delta})$ for all $0\leq i \leq d$.
\end{coro}

\setcounter{section}{3}
\section{The \texorpdfstring{$h''$}{h''}-numbers of  \texorpdfstring{$\hat{\Delta}$}{Delta}} \label{sec:h''-and-socles}
In this section we prove the following extension of Theorem~\ref{known-h'-h''}\eqref{theo2.2_3} to manifolds with boundary. 

\begin{theorem} \label{soc-and-h''}
Let $\Delta$ be a $(d-1)$-dimensional, connected, orientable $\field$-homology manifold with boundary and let $\Theta$ be a generic l.s.o.p.~for $\field[\complet{\Delta}]$. Then for all $i<d$,
\begin{enumerate}
\item\label{theo4.1_1} $\dim_\field \big(\Soc \field(\complet{\Delta},\Theta)\big)_i=\binom{d-1}{i-1}\tilde{\beta}_{i-1}(\complet{\Delta})+\binom{d-1}{i}\tilde{\beta}_{i-1}(\Delta)$, and
\item\label{theo4.1_2} $h''_i(\complet{\Delta})=h_i(\complet{\Delta})-\sum_{j=1}^{i}(-1)^{i-j}\left(\binom{d-1}{i-1}\tilde{\beta}_{j-1}(\complet{\Delta})+\binom{d-1}{i}\tilde{\beta}_{j-1}(\Delta)\right)$.
\end{enumerate}
\end{theorem}
\setcounter{theorem}{2}
To prove Theorem~\ref{soc-and-h''}, several lemmas are in order. As in the previous section, we continue to assume that $\Theta$ is a generic l.s.o.p.~for $\field[\complet{\Delta}]$ and that all coefficients of $\theta_1$ are equal to 1.


\begin{lemma} \label{mod-theta1-decompos}
Let $\Delta$ be a $(d-1)$-dimensional $\field$-homology manifold with boundary. Then 
\[
\big(\Soc\field(\complet{\Delta},\Theta)\big)_i\cong \left(\bigoplus_{j=0}^{d-2}\binom{d-1}{j}\big(H^j_\mideal\big(\field[\complet{\Delta}]/\theta_1\field[\complet{\Delta}]\big)\big)_{i-j}\right)\bigoplus (\SB)_{i-(d-1)} \quad \forall i\in\Z,
\]
where $\SB$ is a graded submodule of $\Soc H^{d-1}_\mideal\big(\field[\complet{\Delta}]/\theta_1\field[\complet{\Delta}]\big)$.
Furthermore, for $j\leq d-2$,
\[
\dim_\field \big(H^j_\mideal\big(\field[\complet{\Delta}]/\theta_1\field[\complet{\Delta}]\big)\big)_\ell=
\begin{cases}
\tilde{\beta}_j(\complet{\Delta}) & \mbox{ if $\ell=1$}\\
\tilde{\beta}_{j-1}(\Delta) & \mbox{ if $\ell=0$}\\
0 & \mbox{ otherwise}.
\end{cases}
\]
\end{lemma}
\begin{proof} 
Since $\complet{\Delta}$ has at most one singularity, Lemma~4.3(2) of~\cite{Novik-Swartz-12} implies that $\field[\complet{\Delta}]/\theta_1\field[\complet{\Delta}]$ is a Buchsbaum $A$-module of Krull dimension $d-1$. The first part of the statement then follows from~\cite[Theorem~2.2]{Novik-Swartz-09:Socles}, while the second part follows from~\cite[Lemma~4.3(1) and Theorem~4.7]{Novik-Swartz-12}.
\end{proof}

We now turn our attention to the submodule $\SB$ of $\Soc H^{d-1}_\mideal\big(\field[\complet{\Delta}]/\theta_1\field[\complet{\Delta}]\big)$.
\begin{prop} \label{SB}
Let $\Delta$ be a $(d-1)$-dimensional, connected, orientable $\field$-homology manifold with boundary. Then, for all $\ell\leq -1$, $\big(\Soc H^{d-1}_\mideal\big(\field[\complet{\Delta}]/\theta_1\field[\complet{\Delta}]\big)\big)_\ell=0$, and hence $(\SB)_\ell=0$.
\end{prop}
\begin{proof}
Since $\depth \field[\complet{\Delta}]\geq 1$, $\theta_1$ is a non-zero divisor on $\field[\complet{\Delta}]$; in other words, the sequence
\[
0\to \field[\complet{\Delta}](-1) \stackrel{\cdot \theta_1}{\longrightarrow} \field[\complet{\Delta}] \longrightarrow \field[\complet{\Delta}]/\theta_1\field[\complet{\Delta}]\to 0
\]
is exact. (For a graded $A$-module $M$, $M(-1)$ denotes $M$ with grading defined by $M(-1)_\ell=M_{\ell-1}$.) 

The above sequence induces a long exact sequence in local cohomology. In particular, the part
\[
\begin{aligned}
H^{d-1}_\mideal\big(\field[\complet{\Delta}]\big)(-1) &\stackrel{\cdot \theta_1}{\longrightarrow} H^{d-1}_\mideal\big(\field[\complet{\Delta}]\big) \longrightarrow H^{d-1}_\mideal\big(\field[\complet{\Delta}]/\theta_1\field[\complet{\Delta}]\big)\longrightarrow H^{d}_\mideal\big(\field[\complet{\Delta}]\big)(-1) 
\\
&\stackrel{\cdot \theta_1}{\longrightarrow} H^{d}_\mideal\big(\field[\complet{\Delta}]\big)
\end{aligned}
\]
is exact. Thus, $H^{d-1}_\mideal\big(\field[\complet{\Delta}]/\theta_1\field[\complet{\Delta}]\big)$, considered as a vector space, is isomorphic to the direct sum of 
\[
\C:=\Coker \big[H^{d-1}_\mideal\big(\field[\complet{\Delta}]\big)(-1) \stackrel{\cdot \theta_1}{\longrightarrow} H^{d-1}_\mideal\big(\field[\complet{\Delta}]\big)\big]\] 
and 
\[\K:=\Ker \big[H^{d}_\mideal\big(\field[\complet{\Delta}]\big)(-1) \stackrel{\cdot \theta_1}{\longrightarrow} H^{d}_\mideal\big(\field[\complet{\Delta}]\big)\big].
\]
 Futhermore, on the $\K$-part of $H^{d-1}_\mideal\big(\field[\complet{\Delta}]/\theta_1\field[\complet{\Delta}]\big)$, the $A$-module structure is induced by the $A$-module structure on $H^{d}_\mideal\big(\field[\complet{\Delta}]\big)$.

Since $\complet{\Delta}$ has (at most) one singular vertex, namely $v_0$, Theorem~\ref{local-cohom-1-singul} implies that for $\ell\leq -1$, the map $\cdot\theta_1: \big(H^{d-1}_\mideal\big(\field[\complet{\Delta}]\big)\big)_{\ell-1}\to \big(H^{d-1}_\mideal\big(\field[\complet{\Delta}]\big)\big)_{\ell}$ is the identity map. Hence its cokernel, $\C_\ell$, vanishes for all $\ell\leq -1$. Therefore, it only remains to show that the socle $(\Soc \K)_\ell$, vanishes for all $\ell\leq -1$. Indeed, by definition of socles,
\begin{align*}
(\Soc \K)_\ell &=\left(\Soc \Ker \big[\cdot \theta_1: H^{d}_\mideal\big(\field[\complet{\Delta}]\big)(-1) \longrightarrow H^{d}_\mideal\big(\field[\complet{\Delta}]\big)\big]\right)_\ell \\
&= \left(\Soc H^{d}_\mideal\big(\field[\complet{\Delta}]\big)(-1)\right)_\ell=\left(\Soc H^{d}_\mideal\big(\field[\complet{\Delta}]\big)\right)_{\ell-1}.
\end{align*}
The following lemma verifies that the latter term vanishes, and thus completes the proof. 
\end{proof}

\begin{lemma} \label{vanishing-socle}
Let $\Delta$ be a $(d-1)$-dimensional, connected, orientable $\field$-homology manifold with boundary. Then, for all $\ell\geq 2$,
$\left(\Soc H^{d}_\mideal\big(\field[\complet{\Delta}]\big)\right)_{-\ell}=0.$
\end{lemma}
\begin{proof} Recall that by Theorem~\ref{Grabe-d},
\begin{equation} \label{eq:H_d}
H^d_{\mideal}(\field[\complet{\Delta}])_{-\ell}=\bigoplus_{\rho\in\M_\ell(\complet{\Delta})} \HH_{\rho}, \quad \mbox{where} \quad \HH_{\rho}=H^{d-1}_{s(\rho)}(\complet{\Delta}).
\end{equation}
Fix $\ell\geq 2$, and let $\rho\in\M_\ell(\complet{\Delta})$. Then either $\rho$ is divisible by $x_v^2$ for some vertex $v$ of $\complet{\Delta}$ (possibly $v_0$) or $\rho$ is a squarefree monomial whose support has size at least two: $s(\rho)\supseteq\{v, w\}$. In the former case, by Theorem~\ref{Grabe-d}, the multiplication map $\cdot x_v: \HH_{\rho} \to \HH_{\rho/x_v}$ is the identity map, and so no non-zero element of $\HH_{\rho}$ is in the socle. In the latter case, at least one of $v,w$ is not $v_0$. Assume without loss of generality that $w\neq v_0$, and consider the map $\cdot x_v: \HH_{\rho} \to \HH_{\rho/x_v}$, which by Theorem~\ref{Grabe-d} is simply $\iota^\ast: H^{d-1}_{s(\rho)}(\complet{\Delta}) \to H^{d-1}_{s(\rho/x_v)}(\complet{\Delta})$. We will show that this map is an isomorphism, and hence that no non-zero element of $\HH_{\rho}$ is in the socle in this case as well.

Our argument is similar to the one used in the proof of~\cite[Theorem~2.1]{Novik-Swartz-09:Gorenstein}. Denote by $\|\complet{\Delta}\|$ the geometric realization of $\complet{\Delta}$, and by $b(\rho)$ and $b(\rho/x_v)$ the barycenters of realizations of faces $s(\rho)$ and $s(\rho/x_v)$, respectively. Consider the following commutative diagram, where the maps $f^\ast$ and $j^\ast$ are induced by inclusion:
\[
\begin{CD}
\tilde{H}^{d-1}\big(\|\complet{\Delta}\|\big) @>(j^\ast)^{-1}>> 
\tilde{H}^{d-1}\big(\|\complet{\Delta}\|, \|\complet{\Delta}\| - b(\rho/x_v)\big)
 @>f^\ast>> \tilde{H}^{d-1}\big(\complet{\Delta}, \cost_{\complet{\Delta}}s(\rho/x_v)\big) \\
 @| @. @A\iota^\ast AA \\
 \tilde{H}^{d-1}\big(\|\complet{\Delta}\|\big) @>(j^\ast)^{-1}>> 
\tilde{H}^{d-1}\big(\|\complet{\Delta}\|, \|\complet{\Delta}\| - b(\rho)\big) 
@>f^\ast>> \tilde{H}^{d-1}\big(\complet{\Delta}, \cost_{\complet{\Delta}}s(\rho)\big). 
\end{CD}
\]
The two maps $f^\ast$ are isomorphisms by the usual deformation retractions. Since $w\neq v_0$ and $w\in s(\rho/x_v) \subset s(\rho)$, the links $\lk_{\complet{\Delta}} s(\rho)$ and $\lk_{\complet{\Delta}} s(\rho/x_v)$ are $\field$-homology spheres, so the four $\field$-spaces on the right and in the middle of the diagram are $1$-dimensional. Furthermore, since $\Delta$ is connected and orientable, the $\field$-spaces on the left of the diagram are $1$-dimensional and the two $j^\ast$-maps are isomorphisms, so that $(j^\ast)^{-1}$-maps are well-defined and are isomorphisms as well. This implies that $\iota^\ast$ is an isomorphism and completes the proof.
\end{proof}

We are now ready to prove Theorem~\ref{soc-and-h''}.

\begin{proof}[Proof of Theorem~\ref{soc-and-h''}]
We prove both parts simulataneously. If $d=2$, then $\complet{\Delta}$ is a circle, in which case the statement is known. So assume $d\geq 3$. Lemma~\ref{mod-theta1-decompos} and Proposition~\ref{SB} imply that the formula for the dimension of the socle holds for all $i\leq d-2$. Together with Theorem~\ref{known-h'-h''}\eqref{theo2.2_4} and Definition~\ref{def-h'-h''}, this also implies that the formula for $h''_i(\complet{\Delta})$ holds for all $i\leq d-2$. Thus, it only remains to show that the theorem holds for $i=d-1$. Since by Corollary~\ref{symmetry}, the $h''$-numbers of $\complet{\Delta}$ are symmetric, to complete the proof of both parts, it suffices to check that the proposed expression for $h''_{d-1}(\complet{\Delta})$ is equal to $h''_1(\complet{\Delta})=h_1(\complet{\Delta})$.

Let $\tilde{\chi}$ denote the reduced Euler characteristic. Note that since $\tilde{\beta}_{d-1}(\complet{\Delta})=1$ and $\tilde{\beta}_{d-1}(\Delta)=0$, the proposed expression for $h''_{d-1}(\complet{\Delta})$, 
\[
h_{d-1}(\complet{\Delta})-\sum_{j=1}^{d-1}(-1)^{d-j-1}\left[(d-1)\tilde{\beta}_{j-1}(\complet{\Delta})+\tilde{\beta}_{j-1}(\Delta)\right],
\]
 can be rewritten as
\[
h_{d-1}(\complet{\Delta})- (d-1)\left(1+(-1)^d\tilde{\chi}(\complet{\Delta})\right)-(-1)^d\tilde{\chi}(\Delta).
\]
Thus to complete the proof, we only need to verify that 
\[
h_{d-1}(\complet{\Delta})=h_{1}(\complet{\Delta})+(d-1)\left(1+(-1)^d\tilde{\chi}(\complet{\Delta})\right)+(-1)^d\tilde{\chi}(\Delta).
\]

To do so, observe that for all $i$,
\[
f_i(\complet{\Delta})=f_i(\Delta)+f_i(\st_{\complet{\Delta}} v_0)- f_i(\partial\Delta).
\]
This, together with the fact that vertex stars are contractible, implies that
\begin{equation} \label{eq:chi}
\tilde{\chi}(\partial\Delta)=\tilde{\chi}(\Delta)+\tilde{\chi}(\st_{\complet{\Delta}} v_0)-\tilde{\chi}(\complet{\Delta})=
\tilde{\chi}(\Delta)-\tilde{\chi}(\complet{\Delta}).
\end{equation}
Finally, according to~\cite[Theorem~3.1]{Novik-Swartz-12},
\begin{align*}
 h_{d-1}(\complet{\Delta})&\ \ = \ \ h_{1}(\complet{\Delta})+d\left(1+(-1)^d\tilde{\chi}(\complet{\Delta})\right)-\left(1+(-1)^{d-1}\tilde{\chi}(\partial\Delta)\right)\\
& \ \ \mathop{=}^{\hbox to 0pt {\hss\scriptsize by~\eqref{eq:chi}\hss}} \ \ 
h_{1}(\complet{\Delta})+(d-1)+ d(-1)^d \tilde{\chi}(\complet{\Delta})+(-1)^d\left(\tilde{\chi}(\Delta)-\tilde{\chi}(\complet{\Delta})\right)\\
&\ \ = \ \ h_{1}(\complet{\Delta})+(d-1)\left(1+(-1)^d\tilde{\chi}(\complet{\Delta})\right)+(-1)^d\tilde{\chi}(\Delta).
\end{align*}
The result follows.
\end{proof}
\begin{thebibliography}{99}
\bibitem[2]{Athanasiadis-11} Athanasiadis, Christos A. Some combinatorial properties of flag simplicial pseudomanifolds and spheres, Ark. Mat., Volume 49 (2011) no. 1, pp. 17-29.
\bibitem[3]{Babson-Novik} Babson, Eric; Novik, Isabella Face numbers and nongeneric initial ideals, Electron. J. Combin., Volume 11 (2004/06) no. 2, Paper no. Research Paper 25, 23 pages.
\bibitem[12]{Grabe-84} Gräbe, Hans-Gert The canonical module of a Stanley–Reisner ring, J. Algebra, Volume 86 (1984) no. 1, pp. 272-281.
\bibitem[19]{Munkres-book} Munkres, James R. Elements of algebraic topology, Addison-Wesley Publishing Company, Menlo Park, CA, 1984, ix+454 pages.
\bibitem[28]{Novik-Swartz-09:Gorenstein} Novik, Isabella; Swartz, Ed Gorenstein rings through face rings of manifolds, Compos. Math., Volume 145 (2009) no. 4, pp. 993-1000.
\bibitem[29]{Novik-Swartz-09:Socles} Novik, Isabella; Swartz, Ed Socles of Buchsbaum modules, complexes and posets, Adv. Math., Volume 222 (2009) no. 6, pp. 2059-2084.
\bibitem[30]{Novik-Swartz-12} Novik, Isabella; Swartz, Ed Face numbers of pseudomanifolds with isolated singularities, Math. Scand., Volume 110 (2012) no. 2, pp. 198-222.
\bibitem[32]{Schenzel-81} Schenzel, Peter On the number of faces of simplicial complexes and the purity of Frobenius, Math. Z., Volume 178 (1981) no. 1, pp. 125-142.
\bibitem[33]{Stanley-75} Stanley, Richard P. The upper bound conjecture and Cohen–Macaulay rings, Studies in Appl. Math., Volume 54 (1975) no. 2, pp. 135-142.
\bibitem[35]{Stanley-96} Stanley, Richard P. Combinatorics and commutative algebra, Progress in Mathematics, 41, Birkhäuser Boston, Inc., Boston, MA, 1996, x+164 pages.
\bibitem[36]{Swartz-09} Swartz, Ed Face enumeration—from spheres to manifolds, J. Eur. Math. Soc. (JEMS), Volume 11 (2009) no. 3, pp. 449-485.
\bibitem[38]{TayWhiteley} Tay, Tiong-Seng; White, Neil; Whiteley, Walter Skeletal rigidity of simplicial complexes. I, European J. Combin., Volume 16 (1995) no. 4, pp. 381-403.
\end{thebibliography}
\end{document}
